\documentclass[10pt,a4paper]{amsart}
\usepackage[margin=19mm]{geometry}
\usepackage{amsmath,amssymb,mathtools}
\usepackage[scr=rsfs,cal=euler]{mathalfa}
\usepackage{xfrac}
\usepackage[unicode,hidelinks,bookmarks=false]{hyperref}
\newcommand{\ie}{i.e.\ }
\newcommand{\cf}{cf.\ }
\newcommand{\resp}{respectively\ }
\newcommand{\Subsection}{\subsection}
\providecommand{\nobreakdash}{\nobreak\hbox{-}}
\setlength{\emergencystretch}{2em}
\multlinegap0pt
\usepackage{amssymb}
\usepackage{tikz}
\newtheorem{lemma}{Lemma}
\newtheorem{theorem}{Theorem}
\title{\vspace{-2.2em}Finite Configurations Cannot Generate a Constant Trace in Rule 30}
\author{David L. Condrey}
\date{Source: arXiv:2609.09431v1 (CC BY 4.0)}
\begin{document}
\maketitle

\noindent\emph{Source and licence.} \vspace{-2.2em}Finite Configurations Cannot Generate a Constant Trace in Rule 30. Version arXiv:2609.09431v1 (CC BY 4.0), distributed by arXiv under the Creative Commons Attribution 4.0 International licence, http://creativecommons.org/licenses/by/4.0/, as recorded in the arXiv licence metadata for this exact version and shown on the abstract page https://arxiv.org/abs/2609.09431v1. Full licence notice: \texttt{licenses/arxiv-2609.09431-CC-BY-4.0.txt}.\par\smallskip

\noindent\textit{Editorial scope.} Counted unit: the article's theorem thm:fiber (source0.tex lines 675--692). The article's complete original proof of it is delivered with it (source0.tex lines 693--723, one \texttt{proof} environment of 31 lines). No mathematics is written, repaired or re-derived: the statement and the proof are the article's own lines, in the article's own notation, and no retained line is substituted or re-flowed.  Retained uncounted as the source-local context the proof needs: lines 656--659.  Nothing was omitted from any retained span, so the package carries no omission record.  No retained line contains a citation, so the wrapper carries no bibliography and no bibliography entry.  No retained line contains a figure environment, an image-inclusion command or drawing code, so no image is delivered and the package declares no asset dependency.  Editorial formatting only, in addition: an AMS \texttt{amsart} preamble that restates the article's own declarations and macros used by the retained lines, namely the environments \texttt{lemma}, \texttt{theorem} on separate counters, together with the standard packages those lines need.  The retained environments print in this document as Lemma 1, Theorem 1 in order of appearance, whereas the article prints the same items with the numbering of its own full document.  The complete native source of the article as distributed in the arXiv e-print for this exact version is bundled unchanged as \texttt{sources/209813/source0.tex}.
\begin{lemma}[Triangular uniqueness]\label{lem:unique}
If two configurations agree at every coordinate $i\geq0$ and have the same
central trace, then they are equal.
\end{lemma}

\begin{theorem}[Zero-trace fiber]\label{thm:fiber}
For a prescribed right half with $R_0=0$, the zero-trace-compatible left half
is unique.  Suppose the positive half is nonzero, and let
$m=\min\{j\geq1:R_j=1\}$.  The forced left half is then the following, with
$R_j$ for $j>m$ arbitrary.

\begin{equation}\label{eq:pattern}
 L_j=0\ (j<m),\qquad L_m=1,\qquad L_j=j\bmod2\ (j>m)
\end{equation}

If instead the positive half is zero, the unique row is zero.  The same
classification in prefix-OR form is the following.

\begin{align}
 L_{2k+1}&=\bigvee_{j=1}^{2k+1}R_j &&(k\geq0) \label{eq:odd}\\[-2pt]
 L_{2k}&=R_{2k}\land\neg\!\bigvee_{j=1}^{2k-1}R_j &&(k\geq1) \label{eq:even}
\end{align}
\end{theorem}

\begin{proof}
For $m\geq1$, let $\mathcal C_m$ consist of rows satisfying these conditions.
One Rule 30 step gives the three identities below.

\begin{equation}\label{eq:updates}
\begin{aligned}
 x'_0&=L_1\oplus R_1\\
 R'_j&=R_{j-1}\oplus(R_j\lor R_{j+1})\\
 L'_j&=L_{j+1}\oplus(L_j\lor L_{j-1})
\end{aligned}
\end{equation}

Suppose $m>1$.  These identities then give the following.

\begin{equation}\label{eq:step}
 x'_0=0,\quad R'_j=L'_j=0\ (j<m-1),\quad R'_{m-1}=L'_{m-1}=1,
 \quad L'_m=m\bmod2
\end{equation}

For $j>m$, including $j=m+1$ because $L_m=1$, the OR in the last update is
one, so $L'_j=(1-(j\bmod2))\oplus1=j\bmod2$.  Thus
$F(\mathcal C_m)\subseteq\mathcal C_{m-1}$.  For $m=1$, the alternating left
half is fixed, $x'_0=1\oplus1=0$, and
$R'_1=0\oplus(1\lor R_2)=1$ independently of the right tail; hence
$F(\mathcal C_1)\subseteq\mathcal C_1$.  Every $\mathcal C_m$ therefore
realizes the zero trace.

A nonzero right half selects a realization by its first one, and
Lemma~\ref{lem:unique} makes it unique; a zero right half forces a zero left
half by the same lemma against the zero row.
\end{proof}

\end{document}
